\subsection{凯莱代数的构造。四元数}

设 \((\mathbf{E}, s)\) 是 A 上的一个凯莱代数（对它我们将使用第 4 号的记号），并设 \(\gamma \in \mathbf{A}\) 。设 F 是 A 上的这样的代数：其底模为 \(\mathbf{E} \times \mathbf{E}\) ，乘法由下式定义

\[
(x, y) (x ^ {\prime}, y ^ {\prime}) = (x x ^ {\prime} + \gamma \bar {y} ^ {\prime} y, y \bar {x} ^ {\prime} + y ^ {\prime} x);\tag{22}
\]

显然 \((e,0)\) 是 F 的单位元，且 \(E \times \{0\}\) 是 F 的一个与 E 同构的子代数；在下文中我们把它与 E 等同，于是 \(x \in E\) 与 \((x,0)\) 等同，特别地，e 与 F 的单位元等同。

设 \(t\) 是 \(\mathbf{F}\) 的由下式定义的置换

\begin{enumerate}
\def\labelenumi{(\arabic{enumi})}
\setcounter{enumi}{22}
\tightlist
\item
\end{enumerate}

\[
t ((x, y)) = (\bar {x}, - y) \quad (x \in \mathbf {E}, y \in \mathbf {E}).
\]

命题 5. (i) 有序对 \((\mathbf{F}, t)\) 是 \(\mathbf{A}\) 上的一个凯莱代数。

\begin{enumerate}
\def\labelenumi{(\roman{enumi})}
\setcounter{enumi}{1}
\tightlist
\item
  设 \(j = (0, e)\) ，从而 \((x, y) = xe + yj\) （对 \(x \in \mathbf{E}\) ， \(y \in \mathbf{E}\) ）。 \(\mathbf{T}_{\mathbb{F}}\) 与 \(\mathbf{N}_{\mathbb{F}}\) （ \(\mathbf{F}\) 的凯莱迹与范数）由下列公式给出
\end{enumerate}

\[
\mathrm{T} _ {\mathrm{F}} (x e + y j) = \mathrm{T} (x), \quad \mathrm{N} _ {\mathrm{F}} (x e + y j) = \mathrm{N} (x) - \gamma \mathrm{N} (y).\tag{24}
\]

\begin{enumerate}
\def\labelenumi{(\roman{enumi})}
\setcounter{enumi}{2}
\tightlist
\item
  要使 F 为结合的，必要且充分的条件是 E 为结合的且交换的。
\end{enumerate}

对 \((x,y)\in \mathbf{F},\)

\begin{enumerate}
\def\labelenumi{(\arabic{enumi})}
\setcounter{enumi}{24}
\tightlist
\item
\end{enumerate}

\[
\begin{array}{r l r} & & {(x, y) + t ((x, y)) = (x + \bar {x}, 0) = \mathrm{T} (x) e} \\ & & {(x, y) t ((x, y)) = (x, y) (\bar {x}, - y) = (x \bar {x} - \gamma \bar {y} y, y \bar {\bar {x}} - y x)} \\ & & {= (\mathrm{N} (x) - \gamma \mathrm{N} (y), 0) = (\mathrm{N} (x) - \gamma \mathrm{N} (y)) e.} \end{array}\tag{26}
\]

因此要证明 (i) 与 (ii)，只需证明 t 是 F 的一个反自同构。显然 t 是一个 A-线性双射。另一方面，

\[
\begin{array}{r l} t ((x, y) \cdot (x ^ {\prime}, y ^ {\prime})) & = t ((x x ^ {\prime} + \gamma \bar {y} ^ {\prime} y, y \bar {x} ^ {\prime} + y ^ {\prime} x)) = (\bar {x} ^ {\prime} \bar {x} + \gamma \bar {y} y ^ {\prime}, - y \bar {x} - y ^ {\prime} x) \\ & = (\bar {x} ^ {\prime}, - y ^ {\prime}) (\bar {x}, - y) = t ((x ^ {\prime}, y ^ {\prime})) t ((x, y)) \end{array}
\]

从而 \(t\) 是一个反自同构。

还需证明 (iii)。由于 E 被等同于 F 的一个子代数，可以假设 E 是结合的。设 \(u = (x, y)\) , \(u' = (x', y')\) , \(u'' = (x'', y'')\) 为 F 的元素。则

\[
\left\{ \begin{array}{l} (u u ^ {\prime}) u ^ {\prime \prime} = ((x x ^ {\prime} + \gamma \bar {y} ^ {\prime} y) x ^ {\prime \prime} + \gamma \bar {y} ^ {\prime \prime} (\bar {y} x ^ {\prime} + y ^ {\prime} x), \\ \qquad \qquad \qquad \qquad \qquad \qquad \qquad \qquad \qquad \qquad (y \bar {x} ^ {\prime} + y ^ {\prime} x) \bar {x} ^ {\prime \prime} + y ^ {\prime \prime} (x x ^ {\prime} + \gamma \bar {y} ^ {\prime} y)) \\ u (u ^ {\prime} u ^ {\prime \prime}) = (x (x ^ {\prime} x ^ {\prime \prime} + \gamma \bar {y} ^ {\prime \prime} y ^ {\prime}) + (\gamma (x ^ {\prime \prime} \bar {y} ^ {\prime} + \bar {x} ^ {\prime} \bar {y} ^ {\prime \prime}) y, \\ \qquad \qquad \qquad \qquad \qquad \qquad \qquad \qquad \qquad \qquad y (\bar {x} ^ {\prime \prime} \bar {x} ^ {\prime} + \gamma \bar {y} ^ {\prime} y ^ {\prime \prime}) + (y ^ {\prime} \bar {x} ^ {\prime \prime} + y ^ {\prime \prime} x ^ {\prime}) x). \end{array} \right.\tag{27}
\]

检查这些公式可知，E 的交换性蕴含 F 的结合性。反之，若 F 是结合的，则把公式 (27) 应用于 \(y = y' = 0\) , \(x'' = 0\) 与 \(y'' = e\) 便得 \((0, x'x) = (0, xx')\) ，即 \(x'x = xx'\) 对 E 中所有 \(x, x'\) 成立；因此 E 此时是交换的。

还需注意，在上述记号中，当 \(x, y\) 属于 \(\mathbf{E}\) 时，

\begin{enumerate}
\def\labelenumi{(\arabic{enumi})}
\setcounter{enumi}{27}
\tightlist
\item
\end{enumerate}

\[
\begin{array}{r l} y j = j \overline {{y}}, & x (y j) = (y x) j, \quad (x j) y = (x \overline {{y}}) j, \quad (x j) (y j) = \overline {{y}} x e \\ & j ^ {2} = e. \end{array}\tag{29}
\]

凯莱代数 \((\mathbf{F}, t)\) 称为 \((\mathbf{E}, s)\) 的由 \(\gamma\) 所定义的凯莱扩张。

例子。(1) 若取 \(\mathbf{E} = \mathbf{A}\) （从而 \(s = 1_{\mathbf{A}}\) ），则代数 \(\mathbf{F}\) 是一个二次 \(\mathbf{A}\) -代数，带有基 \((e,j)\) ，其中 \(j^2 = \gamma e\) 。

\begin{enumerate}
\def\labelenumi{(\arabic{enumi})}
\setcounter{enumi}{1}
\tightlist
\item
  取 E 为 \((\alpha, \beta)\) 型的二次代数，于是 E 的底模为 \(A^{2}\) ，其典范基的乘法表为 (2)（第 3 号）。取 s 为 E 的共轭（第 3 号，命题 2）。于是对所有 \(\gamma \in A\) ， \((E, s)\) 的由 \(\gamma\) 所定义的凯莱扩张 F 称为 \((\alpha, \beta, \gamma)\) 型的四元数代数，由第 3 号命题 1 与上述命题 5，它是结合的；其底模为 \(A^{4}\) ，且若 \((e, i, j, k)\) 表示 \(A^{4}\) 的典范基，则相应的乘法表由
\end{enumerate}

\[
\left\{ \begin{array}{l l} i ^ {2} = \alpha e + \beta i, & i j = k, \qquad i k = \alpha j + \beta k, \\ j i = \beta j - k, & j ^ {2} = \gamma e, \qquad j k = \beta \gamma e - \gamma i, \\ k i = - \alpha j, & k j = \gamma i, \qquad k ^ {2} = - \alpha \gamma e. \end{array} \right.\tag{30}
\]

进一步，对 \(u = \rho e + \xi i + \eta j + \eta k\) （其中 \(\rho, \xi, \eta, \zeta\) 属于 A），我们有（以 \(\bar{u}\) 代替 \(t(u)\) 来记，并将 A 与 Ae 等同）：

\[
\left\{ \begin{array}{c} u = (\rho + \beta \xi) e - \xi i - \eta j - \zeta k \\ \mathrm{T} _ {\mathrm{F}} (u) = 2 \rho + \beta \xi \\ \mathrm{N} _ {\mathrm{F}} (u) = \rho^ {2} + \beta \rho \xi - \alpha \xi^ {2} - \gamma (\eta^ {2} + \beta \eta \zeta - \alpha \zeta^ {2}). \end{array} \right.\tag{31}
\]

公式 (30) 由 (28) 与 (29) 得出，公式 (31) 由 (23) 与 (24) 得出，并考虑到二次代数 E 的诸公式。

于是，当 \(u, v\) 属于 \(\mathbf{F}\) 时，

\[
\mathrm{N} _ {\mathrm{F}} (u v) = \mathrm{N} _ {\mathrm{F}} (u) \mathrm{N} _ {\mathrm{F}} (v)\tag{32}
\]

因为 \(\mathrm{N}_{\mathbb{F}}(uv)=uv.\overline{uv}=uv(\bar{v}.\bar{u})=u(v\bar{v})\bar{u}=(u\bar{u})(v\bar{v})\) ，这由结合性以及 \(\mathrm{N}_{\mathbb{F}}(u)\) 属于 F 的中心这一事实得出。

与四元数代数同构的 A-代数也称为四元数代数；若这样一个代数的基具有乘法表 (30)，则称该基为 \((\alpha, \beta, \gamma)\) 型的基。作为语言上的滥用，当一个四元数代数具有 \((\alpha, \beta, \gamma)\) 型的基时，就说它是 \((\alpha, \beta, \gamma)\) 型的。

当 \(\beta = 0\) 时，公式 (30) 与 (31) 简化为

\[
\left\{ \begin{array}{l l} i ^ {2} = \alpha e, & \quad i j = k, \qquad i k = \alpha j, \\ j i = - k, & \quad j ^ {2} = \gamma e, \qquad j k = - \gamma i, \\ k i = - \alpha i, & \quad k j = \gamma i, \qquad k ^ {2} = - \alpha \gamma e, \end{array} \right.\tag{33}
\]

和

\[
\left\{ \begin{array} {c}\bar {u} = \rho e - \xi i - \eta j - \zeta k \\ \mathrm{T} _ {\mathrm{F}} (u) = 2 \rho \\ \mathrm{N} _ {\mathrm{F}} (u) = \rho^ {2} - \alpha \xi^ {2} - \gamma \eta^ {2} + \alpha \gamma \zeta^ {2}. \\ \end{array} \right.\tag{34}
\]

于是在上述表达式中， \((\alpha, \beta, \gamma)\) 处处被 \((\alpha, \gamma)\) 替换。 \((\alpha, \gamma)\) 型与 \((\gamma, \alpha)\) 型的四元数代数显然同构。

注意，公式 (32) 表明，当 \(-1 \neq 1\) 在 A 中时，F 不是交换的。

取 A 为实数域 \(\mathbf{R}\) 且 \(\alpha = \gamma = -1\) , \(\beta = 0\) ，相应的代数 \(F\) 称为哈密顿四元数代数，记作 \(\mathbf{H}\) 。若 \(u = \rho e + \xi i + \eta j + \zeta k\) （ \(\rho, \xi, \eta, \zeta\) 属于 \(\mathbf{R}\) ）是一个 \(\neq 0\) 的、属于 \(\mathbf{H}\) 的元素，则公式 \(u\bar{u} = \bar{u} u = \rho^2 + \xi^2 + \eta^2 + \zeta^2\) （参见 (34)）表明 \(N(u) \neq 0\) 在 \(\mathbf{R}\) 中，于是 \(u\) 具有一个逆元 \(u^{-1} = N(u)^{-1}\bar{u}\) （在 \(\mathbf{H}\) 中），因而 \(\mathbf{H}\) 是一个非交换域。

\begin{enumerate}
\def\labelenumi{(\arabic{enumi})}
\setcounter{enumi}{2}
\tightlist
\item
  若将 E 取为四元数代数（参见例 2），则由元素 \(\delta \in A\) 所定义的 E 的凯莱扩张一般是非结合的（命题 5）；它称为 A 上的一个八元数代数（参见附录，第 3 号）。
\end{enumerate}

